\section*{Capítulo \ref{ch:b2:thermodynamic-potentials} --- Potenciales termodinámicos}

\begin{solution}{exo:b2:thermodynamic-potentials:1}
$F = -nRT\ln V + f(T)$, $G = nRT\ln P + g(T)$; $P = -\partial_VF = nRT/V$,
$V = \partial_PG = nRT/P$; $-\partial_TF = nR\ln V - f'(T)$, la entropía
$nC_V\ln T + nR\ln V + \text{const.}$ si $f' = -nC_V\ln T + \dots$
\end{solution}

\begin{solution}{exo:b2:thermodynamic-potentials:2}
$\Delta F = -RT\ln 10 = \qty{-5.7}{kJ}$, $\Delta U = 0$, $\Delta S = \qty{19.1}{J/K}$.
Reversible: \qty{5.7}{kJ}, $S_{\text{c}} = 0$; contra \qty{1}{bar}: \qty{0.9}{kJ},
$S_{\text{c}} = 19.1 - 900/300 = \qty{16}{J/K}$; al vacío: nada, $S_{\text{c}} =
\qty{19.1}{J/K}$.
\end{solution}

\begin{solution}{exo:b2:thermodynamic-potentials:3}
$\Delta g = \Delta h - T\Delta s$ con $\Delta h \approx L$, $\Delta s \approx L/T_{\text{fus}}$.
A \qty{-5}{\celsius}: $g_{\text{liq}} - g_{\text{hielo}} = +\qty{6.1}{kJ/kg}$, el hielo es
estable y el líquido metastable (existe el subenfriamiento); a $+\qty{5}{\celsius}$:
\qty{-6.1}{kJ/kg}, el líquido es estable --- el hielo sobrecalentado no se observa, pues su
superficie funde primero.
\end{solution}

\begin{solution}{exo:b2:thermodynamic-potentials:4}
Fusión $\qty{-1.3e7}{Pa/K}$; vaporización \qty{3.6}{kPa/K}. Con \qty{70}{bar}:
\qty{-0.5}{K}; a \qty{0.7}{bar}: \qty{92}{\celsius}; a \qty{2}{bar}: $\approx
\qty{127}{\celsius}$ por la pendiente local (\qty{120}{\celsius} en realidad: la
pendiente crece).
\end{solution}

\begin{solution}{exo:b2:thermodynamic-potentials:5}
(a) También $(\partial T/\partial V)_S = -(\partial P/\partial S)_V$ y
$(\partial T/\partial P)_S = (\partial V/\partial S)_P$. (b) $0$; $a/v^2$. (c)
$\alpha/\kappa_T = \qty{4.6}{bar/K}$; $T\alpha/\kappa_T - P \approx \qty{1400}{bar}$.
(d) \qty{46}{bar}: la botella revienta.
\end{solution}

\begin{solution}{exo:b2:thermodynamic-potentials:6}
(a) \cref{thm:b2:thermodynamic-potentials:maxwell}. (b) $nR$. (c) Agua,
\qty{29}{J/kg/K} (el $0.7\%$); cobre, \qty{12}{J/kg/K} (el $3\%$). (d) La
diferencia es despreciable para la materia condensada.
\end{solution}

\begin{solution}{exo:b2:thermodynamic-potentials:7}
(a) $-C(L - L_0)$; $0$. (b) $\Delta S = -C(L - L_0)^2/2 = \qty{-5e-3}{J/K}$; $Q =
\qty{-1.5}{J}$; $W = +\qty{1.5}{J}$. (c) $+\qty{0.4}{K}$. (d) $(L - L_0) \propto
1/T$: $-9\%$.
\end{solution}

\begin{solution}{exo:b2:thermodynamic-potentials:8}
(a) La materia va al $\mu$ menor: $\mu_{\text{v}} < \mu_{\text{liq}}$ cuando $p <
p_{\text{sat}}$; humedad $= p/p_{\text{sat}}$. (b) $-RT\ln 0.5 = \qty{1.7}{kJ/mol}$
a favor del vapor. (c) La evaporación solo necesita $p < p_{\text{sat}}$,
no la ebullición; con un $100\%$ los potenciales son iguales. (d) Sobre las superficies
más frías, cuya $p_{\text{sat}}(T_{\text{s}})$ ha caído por debajo de la $p$ del aire
--- las que radian hacia el cielo nocturno.
\end{solution}

\begin{solution}{exo:b2:thermodynamic-potentials:9}
(a) Estándar. (b) $b = RT_{\text{c}}/8P_{\text{c}} = \qty{4.3e-5}{m^3/mol}$, $a =
27R^2T_{\text{c}}^2/64P_{\text{c}} = \qty{0.36}{Pa.m^6/mol^2}$. (c) $\dd g = v\,\dd P$
a $T$ fija y $g$ igual en los dos extremos: $\int v\,\dd P = 0$ alrededor del
lazo. (d) Inestable donde $\partial_vP > 0$; metastable entre medias: el
líquido sobrecalentado (que salta) y el vapor sobresaturado (niebla que
espera núcleos, cámara de niebla).
\end{solution}

\begin{solution}{exo:b2:thermodynamic-potentials:10}
(a) $\dd P/\dd T = L_{\text{m}}P/RT^2$: $\ln(P/P_0) = -(L_{\text{m}}/R)(1/T - 1/T_0)$.
(b) \qty{28}{mbar} frente a \qty{23}{mbar} --- $L$ es mayor a \qty{20}{\celsius}.
(c) \qty{820}{Pa} frente a \qty{611}{Pa}. (d) $\partial(\Delta g/T)/\partial T = -L/T^2$
y $\Delta g = 0$ a lo largo de la curva dan la misma pendiente.
\end{solution}

\begin{solution}{exo:b2:thermodynamic-potentials:11}
(a) Igualdad de $\mu$ a través de la membrana: $\mu^\circ(P + \Pi) - RTx_{\text{s}} =
\mu^\circ(P)$, es decir, $v_{\text{m}}\Pi = RTx_{\text{s}}$, $\Pi = cRT$. (b) \qty{27}{bar}.
(c) $\Pi V = \qty{2.7}{MJ} = \qty{0.75}{kWh}$ por metro cúbico; las plantas reales,
cuatro veces más. (d) El agua fluye hacia el potencial menor: hacia la
célula rica en soluto, y fuera de ella en salmuera.
\end{solution}

\begin{solution}{exo:b2:thermodynamic-potentials:12}
(a) $S = S_0 - k_BL^2/2Na^2$; $f = k_BTL/Na^2$. (b) $k_BT/Na^2 \propto T$. (c)
\qty{1.7e-13}{N}; $1.8 \times 10^{14}$. (d) $E \sim nk_BT \approx \qty{0.4}{MPa}$.
\end{solution}

\begin{solution}{pb:b2:thermodynamic-potentials:1}
\textbf{1.} $\dd U = T\dd S + f\dd L$; $\dd F = -S\dd T + f\dd L$.

\textbf{2.} Schwarz sobre $\dd F$.

\textbf{3.} $f - T\partial_Tf = CT(L - L_0) - CT(L - L_0) = 0$: la energía
no cambia; el tirón es entrópico.

\textbf{4.} $\Delta S = \qty{-5e-3}{J/K}$; $Q = \qty{-1.5}{J}$; $W = +\qty{1.5}{J}$; suma
nula.

\textbf{5.} $+\qty{0.4}{K}$.

\textbf{6.} $L - L_0 \propto 1/T$: \qty{9.1}{cm}, una contracción de \qty{9}{mm}; los
radios calentados se acortan, el centro de masas de la llanta se desplaza y la rueda
gira.

\textbf{7.} $S = S_0 - k_BL^2/2Na^2$; $f = k_BTL/Na^2$; la penalización de entropía se
paga en $k_BT$.

\textbf{8.} \qty{1.7e-13}{N}; $1.8 \times 10^{14}$ cadenas.

\textbf{9.} Su fuerza es energética ($\partial_LU \ne 0$); calentarlo ablanda
los enlaces, así que $f$ cae con $T$.

\textbf{10.} Pendientes $-s$, con $s_{\text{gas}} > s_{\text{liq}} > s_{\text{hielo}}$; la
línea más baja es la fase estable; los cruces están en $T_{\text{fus}}$ y $T_{\text{vap}}$.

\textbf{11.} \cref{thm:b2:thermodynamic-potentials:coexistence}.

\textbf{12.} $\qty{-1.3e7}{Pa/K}$; negativa porque $v_{\text{liq}} < v_{\text{hielo}}$.

\textbf{13.} \qty{70}{bar}: \qty{-0.5}{K} --- no es así como se patina a
\qty{-10}{\celsius}; lo hacen el rozamiento y la capa superficial cuasilíquida.

\textbf{14.} \qty{270}{bar}, \qty{-2}{K}: con el flujo geotérmico, la
base funde y el casquete desliza.

\textbf{15.} \qty{3.6}{kPa/K}; \qty{92}{\celsius}; unos \qty{127}{\celsius}
(\qty{120}{} en realidad).

\textbf{16.} \qty{28}{mbar} (\qty{23}{}), \qty{820}{Pa} (\qty{611}{}): $L$ crece cuando $T$
baja; errores del $20$ al $30\%$.

\textbf{17.} El mismo $\Delta v$, así que el cociente de pendientes vale $L_{\text{sub}}/L_{\text{vap}}
= 1.15$: la curva de sublimación es más empinada --- un pico en el punto triple.

\textbf{18.} \qty{12}{kJ/kg}; el hielo debe nuclear, y un núcleo pequeño cuesta
más energía de superficie de la que gana: las gotitas puras esperan, hasta
\qty{-40}{\celsius}.

\textbf{19.} $\Delta v \to 0$, $L \to 0$, con cociente finito: la curva se acaba;
más allá, un solo fluido que pasa de forma continua de parecerse al líquido a parecerse al gas.

\textbf{20.} $\dd g = v\,\dd P$ a $T$ fija; $g$ igual en los dos extremos de
la horizontal: $\int v\,\dd P = 0$, áreas iguales.

\textbf{21.} La rama del líquido sobrecalentado: ningún núcleo de burbuja en una
taza lisa; una perturbación lo hace hervir de golpe.

\textbf{22.} Congélese, bómbese por debajo de \qty{611}{Pa} y caliéntese con suavidad: el camino
cruza la curva de sublimación y nunca la de fusión; la estructura
se conserva.

\textbf{23.} $\mu_{\text{v}}(p) < \mu_{\text{liq}}$ cuando $p < p_{\text{sat}}$; iguales en
la saturación.

\textbf{24.} \qty{120}{\celsius} en vez de \qty{100}{}: cuatro veces más deprisa.

\textbf{25.} $S$ para el aislado, $F$ para $T,V$ fijos y $G$ para $T,P$ fijos;
$-\Delta F$ y $-\Delta G$ son el trabajo recuperable; las relaciones de Maxwell convierten
derivadas de $S$ en la ecuación de estado.
\end{solution}
